\section{Frontier Replay: reopen the reason an opportunity failed}
\label{sec:replay}
\subsection{A release is a trigger to investigate}
The three operational orders of GENESIS are to see the frontier, see its consequences, and improve the machinery that discovers those consequences \citep[\S3]{boucher_genesis}. Constructive challenge supplies a concrete question at each order. What changed in usable capability? Which previously binding condition might that change affect? How should the discovery process itself change when the answer is known?

A capability delta should identify a specific difference: a permitted task now succeeds more reliably, a measurement becomes affordable, a tool gains a necessary interface, or a complete construction requires fewer resources. A product announcement or model label is insufficient. Verify access and test the relevant behavior. The same advance may also strengthen competitors or customer self-service, so a technical gain can coincide with a smaller commercial margin.

A compact Frontier Replay record retains the valuable intended outcome; the best complete route previously considered; the decisive unresolved condition or supported failure; the evidence and source versions; the missing means; and what new observation would justify reopening the decision. It stores the \emph{reason} for the outcome rather than a large transcript of the attempt. Repeated failures should consolidate into applicability rules, tests or useful instruments.

\subsection{Two searches must meet}
Forward search begins with a verified capability change and asks what it could enable. Backward search begins with consequential unmet outcomes and asks what means are missing. Their meeting point is a candidate construction whose dependencies can be checked. Joint dependencies matter: an inexpensive diagnostic might be useless without an acceptable sharing arrangement, while a sharing arrangement is useless without a valid diagnostic.

\begin{figure}[htbp]
\centering
\begin{tikzpicture}[node distance=.26in and .22in,every node/.style={font=\small,align=center},box/.style={draw=Rule,rounded corners=3pt,fill=Ivory,text width=2.0in,minimum height=.62in,inner sep=7pt},arr/.style={-{Latex[length=2.4mm]},line width=.85pt,color=Sea}]
\node[box] (front) {\textbf{Changed usable capability}\\What became accessible?};
\node[box,right=1.05in of front] (need) {\textbf{Consequential blocked outcome}\\What still prevents it?};
\node[box,below right=.45in and -.1in of front,text width=2.8in] (move) {\textbf{Candidate missing move}\\An executable means and a decisive test};
\node[box,below=.35in of move,text width=2.8in] (whole) {\textbf{Complete productive arrangement}\\Acceptance, participants, resources and timing};
\draw[arr] (front.south)--(move.north west);\draw[arr] (need.south)--(move.north east);\draw[arr] (move)--(whole);
\end{tikzpicture}
\caption{Frontier Replay connects a measured capability delta to a specific obstacle. An attractive dependency diagram is a hypothesis until the interfaces, construction and complete outcome are tested.}
\end{figure}

Replay cannot be the entire discovery system. A backlog encodes yesterday's imagination. A possibility challenger must also search for valuable outcomes that were never represented in it, including new beneficiaries, ways to combine capabilities, and ways to remove a presumed task altogether. Allocate investigation by expected information and construction value; do not impose an arbitrary fixed fraction on every mission.

\subsection{Use the strongest old and new alternatives}
For a claim that a frontier change enabled the new route, compare old and new capability access while holding the intended outcome, permissions and evaluation contract fixed. Allow the old-capability comparator to use its strongest feasible alternative. Otherwise the study shows only that one old implementation fails, not that the new capability removed a genuine barrier.

For an institutional contribution, compare the proposed machinery with direct use of the same frontier model and tools. For improved discovery, compare old and new discovery procedures with the same starting information and stock of capabilities. These are separate questions. A general form is
\begin{equation}
Y_{m,p,s}=\text{complete outcome under model }m,\ \text{procedure }p,\ \text{stock }s.
\end{equation}
A crossed design can vary old/new $m$, old/new $p$, and old/new $s$. At fixed model and stock, $Y_{m,p_1,s}-Y_{m,p_0,s}$ estimates a procedure effect under the experiment. At fixed procedure and model, a stock difference concerns inheritance. Interactions are informative and should be reported rather than assigned arbitrarily to the most attractive story \citep[\S12]{boucher_genesis}.

A useful frontier-adaptive institution may win through integration, acceptance or opportunity selection even when its model components are publicly available. That contribution has to pay for its overhead. If direct use of an existing agent is the strongest complete route, adopting it is operationally successful and does not establish superiority of the added discovery mechanism.

\section{Invent the complete productive arrangement}
\label{sec:economy}
\subsection{The criticism may be economic or institutional}
A technical design can be correct yet lack a willing beneficiary. A set of beneficiaries can have positive aggregate surplus while one required participant rationally refuses. A profitable projection can fail before the first receipt because the construction must be paid for earlier. A new capability can eliminate the very recurring operation on which a proposed revenue mechanism depends. These are distinct constraints, developed in \emph{Invent the Economy} \citep[\S2--9]{boucher_economy}.

Market design has long treated implementation details, experimentation and computation as complements to abstract incentives \citep{marketdesign}. The present architecture applies the same discipline to a proposed inventive enterprise. The output is a complete arrangement, not a persuasive description of a business. It names the useful outcome, technical means, actual or hypothetical counterparties, responsibilities, information, risk allocation, rights, payment conditions, failure handling and time-dependent resources.

A possibility challenger asks a constructive economic question: which capability \emph{and} voluntary arrangement could make this valuable work worth undertaking? It does not fabricate buyers or turn assumed willingness to pay into acceptance. Public evidence can motivate an offer; only a counterparty's authorized action can accept it.

\subsection{Running example: a shared qualification workbench}
Consider a deliberately hypothetical case. Several small teams need trustworthy qualification evidence for short production runs. Each team's standalone effort appears too costly. An initial proposal uses an AI-generated analysis to replace much of the testing. The central question is whether the required outcome can actually be established for the relevant operating conditions.

\textbf{First challenge: the evidence does not establish the property.} The critic identifies a case in which the proposed inputs admit two behaviors with different acceptance outcomes. An adjudicator reproduces that ambiguity. A more confident report cannot resolve it. Responses include obtaining a discriminating measurement, changing the production method, procuring an existing qualification service, or rejecting the route.

\textbf{Second challenge: the analysis omitted a shared instrument.} A possibility challenger hypothesizes that the teams' needs have a common experimental structure. A programmable workbench plus task-specific adaptations might reduce repeated construction. The constructor must establish which features are shared, build the interfaces, and test exceptions. If a standard laboratory already offers the same accepted outcome at a better complete cost, that route wins. The paper does not claim that this field or its market is new.

\textbf{Third challenge: shared technology alone does not secure participation.} Some teams may reject common terms, need different evidence, or face a cheaper outside option. A proposed arrangement must allocate preparation costs, usage, rights, support and failure exposure in ways each participant accepts. Conditional commitments and staged delivery are candidate mechanisms, not assumed contracts.

\textbf{Fourth challenge: the candidate cannot fund its own first step.} Even positive contribution after delivery does not pay for construction before receipts. The complete cash schedule may favor a smaller cohort, a different sequence, authorized deposits, supported financing, or non-entry. The frontier may change one of those choices by lowering the cost of a necessary means. It does not repeal cash conservation.

\begin{table}[htbp]
\small
\begin{tabularx}{\textwidth}{P{1.45in}YY}
\toprule
\textbf{Candidate missing move} & \textbf{What must be constructed or established} & \textbf{What could defeat it}\\
\midrule
A discriminating test & A measurement that separates consequential behaviors under relevant conditions. & Unmodeled behavior, calibration error, or a cheaper adequate test.\\
A shared workbench & A valid common core, tested adaptations and capacity allocation. & Exception costs, common failure exposure, or a superior existing service.\\
A cooperative arrangement & Specific voluntary terms, rights and responsibilities with adequate outside-option checks. & A participant refuses; a rival supplies a better complete offer.\\
A feasible sequence & Resources arrive before obligations fall due; failure branches remain fundable. & Late receipts, unsupported financing, refunds or protected obligations.\\
A successor instrument & A retained construction useful on a different task after revalidation. & The task changes outside scope, maintenance dominates, or fresh construction wins.\\
\bottomrule
\end{tabularx}
\caption{A hypothetical industrial illustration, not an executed deployment. The aim is to connect a challenge to productive means and a falsifiable consequence.}
\end{table}

\subsection{Participation, profit and cash are different tests}
For participant $i$, let $U_i(P)$ be expected utility from arrangement $P$ and $U_i^{\mathrm{out}}$ the best relevant outside option. A necessary modeled participation condition is
\begin{equation}
U_i(P)\geq U_i^{\mathrm{out}}\quad\text{for every required participant }i.
\end{equation}
This does not prove actual consent, truthful reporting or equilibrium. If claims depend on strategic incentives, identify feasible deviations, information and enforcement assumptions, then test or derive the specific incentive claim. New capabilities can create new deviations, including combinations of individually permitted actions.

For a finite project and declared horizon, define expected project cash contribution using all lifecycle receipts and outflows, with any residual obligations covered by explicitly specified provisions:
\begin{equation}
\Pi_{\mathrm{cash}}=\E[\text{project cash receipts}-\text{project cash outflows}-\text{residual provisions}],
\end{equation}
This is a project cash measure, not accounting profit, available cash at every date, or social welfare. State the treatment of financing inflows, principal repayment, taxes and residual value; exclude financing inflows from operating receipts and report funding separately. Set provisions only for obligations not already paid or included elsewhere, and reconcile them when settled. Include acquisition, evidence access, generation, criticism, adjudication, failed exploration, construction, production, acceptance, support, maintenance and financing where applicable. Do not count a cost twice by including it in both a branch payoff and a global charge. Funding is not customer revenue. Unknown actual billing is not zero.

For a known ordered cash path with net events $x_1,\ldots,x_n$ before initial unrestricted cash, the minimum starting cash that keeps every prefix nonnegative is
\begin{equation}
B_{\min}=\max\left\{0,\ -\min_{0\leq k\leq n}\sum_{j=1}^{k}x_j\right\}.
\label{eq:cash}
\end{equation}
This elementary deterministic bound assumes the event path is fixed and that reserved obligations have already been accounted for. A contingent project needs a policy and liquidity checks on the relevant branches. A hoped-for receipt is not an observed receipt; future upside does not become available funds merely because it appears in a plan.

As a hand-worked illustration, suppose a project spends 60 arbitrary currency units on a shared core, spends another 30 on adaptations, and then receives 120 after accepted delivery. Its final contribution is 30, but Equation~\eqref{eq:cash} requires 90 initially. Starting with 50 fails. An assumed, authorized 40-unit deposit could make this \emph{success path} fundable, leaving an 80-unit final receipt and the same total contribution. It also creates delivery or refund obligations that must be priced and funded on failure paths. This is arithmetic about supplied terms, not evidence of buyer agreement or new empirical market creation.

\subsection{Joint construction may be the actual object of invention}
A useful transformation may require several mutually supporting changes. More accurate testing can lower counterparties' uncertainty; agreed participation can justify constructing the test; a feasible stage can establish the first required resource. The resulting plan must have a real bootstrap. Drawing arrows in a cycle does not finance or construct its first step.

Represent a candidate dependency as a conditional operation: specified inputs, resources and permissions produce an output and evidence under named assumptions. Combine operations only when interface meaning, capacity, timing and obligations are compatible. A hypothesis of shared structure becomes reusable only after the actual common construction works and the exceptions are understood.

This perspective permits an ambitious outcome without manufacturing dependence on a proprietary apparatus or settlement token. The customer benefits from accepted productive capability. The institution should earn from useful construction and service that survives alternatives. It should not preserve unnecessary steps merely because those steps are billable. If a simpler ordinary route delivers more value, it belongs in the institution's repertoire.

\subsection{Value creation does not guarantee value capture}
The frontier moves for everyone. A construction that creates substantial value may be copied, commoditized, or bypassed. A buyer may perform the work directly. A partner may capture the margin because it supplies a scarce complement. Evaluate price, defensibility and timing against those responses instead of treating a productivity improvement as guaranteed profit.

The practical commercial target is a complete, repeatable offer with identifiable acceptance, credible obligations, real counterparties and supported economics. The stronger research target is to discover such arrangements more effectively than capable alternatives. These objectives can reinforce one another, but neither is established by naming a promising industry or publishing this paper.
